\addtotoc{Хураангуй}

\vspace{10mm}

\begin{center}
\textbf{\large Хураангуй}
\end{center}

\sloppy

ABU Robocon 2026 тэмцээний "Kung Fu Quest" сэдвийн хүрээнд R2 бүрэн автомат робот нь Мейхуа ой (Meihua Forest)-гоос зөв Күнг Фү судрыг таньж, Тик-Так-Тое (Tic-Tac-Toe) тавцанд хүний оролцоогүйгээр судрыг байрлуулах шаардлагатай. Уг шаардлагын үндсэн дээр объект таних (object detection) болон шийдвэр гаргах алгоритмыг нэгтгэж хөгжүүлнэ.

Энэхүү судалгааны ажил нь R2 роботын программчлалын үндсэн хоёр алгоритмыг боловсруулахад чиглэнэ. Объект таних алгоритмд YOLOv12-ийг сонгон хэрэгжүүлж, шийдвэр гаргах алгоритмд Finite State Machine(Төгсгөлөг Төлөвт машин)-ыг ашигласан. Судалгааны туршилт болон хөгжүүлэлтийг Unity симуляцийн орчинд гүйцэтгэсэн бөгөөд YOLOv12-ын машин сургалтад шаардлагатай датасетийг мөн тус орчныг ашиглаж бэлтгэсэн. Хөгжүүлсэн алгоритмуудын ажиллагааг симуляцийн орчинд туршиж, үр дүнг танилцуулна.

Объект таних алгоритм сургалтын үр дүн: симуляцийн орчинд туршилт явуулахад бүрэн үзэгдэж буй объектуудыг 90\%-ээс өндөр хувийн итгэлтэйгээр таньж байсан бол тал хэсэг нь халхлагдсан объектуудыг 70\%-аас өндөр хувийн итгэлтэйгээр таньж чадаж байсан. Мөн туршилтын үед камерийн байрлалыг төрөл бүрийн өнцөгөөс туршихад ойролцоогоор 12 объектийн 11-ийг нь зөв таньж байсан.

\vspace{3mm}
\textit{Түлхүүр үгс: машин сургалт, объект таних алгоритм, шийдвэр гаргах алгоритм, Finite State Machine, YOLOv12, Unity}

\vspace{10mm}

\begin{center}
\textbf{\large Abstract}
\end{center}

\begin{english}
Under the theme of ABU Robocon 2026, Kung Fu Quest, the R2 fully automated robot needs to identify the correct Kung Fu scriptures from Meihua Forest and place the scriptures on the Tic-Tac-Toe board without human intervention. Based on these requirements, object recognition and decision-making algorithms will be integrated and developed.

This research work focuses on the development of two basic algorithms for R2 robot programming. YOLOv12 was selected for the object recognition algorithm and Finite State Machine was used for the decision making algorithm. Research testing and development were performed in the Unity simulation environment, and the datasets required for YOLOv12 machine learning were also prepared using the same environment. The performance of the developed algorithms is tested in a simulation environment and the results are presented.

Object recognition algorithm training results: when tested in a simulated environment, fully visible objects were recognized with a higher than \textbf{90\%} confidence level, while partially obscured objects were recognized with a higher than \textbf{70\%} confidence level. And when testing the camera's position from various angles during testing, it correctly identified \textbf{11 out of 12} objects.

\vspace{3mm}
\textit{Keywords: machine learning, object detection algorithm, decision-making algorithm, Finite State Machine, YOLOv12, Unity}
\end{english}
